% Part: normal-modal-logic
% Chapter: tableaux
% Section: countermodels

\documentclass[../../../include/open-logic-section]{subfiles}

\begin{document}

\olfileid{nml}{tab}{cou}
      
\olsection{له \usetoken{P}{tableau} څخه ضدمدلونه}

\begin{explain}
  د بشپړتيا د قضيې ثبوت يوازې دا نۀ ښيي چې که
  $\Entails !A$ وي، نو $\Proves !A$ وي؛ که
  $\Entails/ !A$ وي، نو د~$!A$ د ضدمدل د جوړولو لاره هم راکوي۔
  په منجمده سرچينه کښې د ناستلزام تر څنګ د~$!A$ پر ځاے
  يوازې $A$ ليکل شوے؛ دلته له اصلي فارمول سره سم~$!A$ راغلے دے۔
  د~$\Log{K}$ په حالت کښې دا لاره د \emph{پرېکړې کړنلاره}
  جوړوي۔ فرض کړئ $\Entails/ !A$ وي۔ د
  \olref[cpl]{prop:complete-tableau} ثبوت د بشپړ
  !!{tableau} د جوړولو طريقه راکوي، او دا طريقه تل درېږي۔
  د~$\Log{K}$ قضيېيزې قاعدې يوازې د لږې پېچلتيا
  مخکښ‌لرونکي !!{formula}s ورزياتوي؛ يعنې هره قضيېيزه قاعده
  پر څانګه د هر نښه لرونکي فارمول
  $\sFmla{S}{!A}[\sigma]$ لپاره يوازې يو ځل پلې کېدے شي۔
  نوي مخکښونه يوازې د
  \iftag{prvBox}{$\TRule{\False}{\Box}$}{}\iftag{notprvBox,notprvDiamond}{}{
    او }\iftag{prvDiamond}{$\TRule{\True}{\Diamond}$}{}
  \iftag{notprvBox,notprvDiamond}{قاعده}{قاعدې} جوړوي؛ دا هم
  يوازې يو ځل پلې کېږي او يو نوے مخکښ جوړوي۔
  \iftag{prvBox}{$\TRule{\True}{\Box}$}{}\iftag{notprvBox,notprvDiamond}{}{
    او }\iftag{prvDiamond}{$\TRule{\False}{\Diamond}$}{}
  \iftag{notprvBox,notprvDiamond}{قاعده}{قاعدې} کېدے شي څو
  ځله پلې شي، خو د هر مخکښ لپاره يوازې يو ځل؛ او نوي
  مخکښونه هم يوازې متناهي ډېر جوړېږي۔ نو له متناهي پړاوونو
  وروسته جوړښت يا تړلې څانګه ورکوي يا بشپړه څانګه۔

  کله چې داسې تابلو جوړه شي چې يوه پرانيستې بشپړه څانګه ولري،
  د \olref[cpl]{thm:tableau-completeness} ثبوت څرګند مدل
  راکوي چې د اصلي مخکښ‌لرونکو !!{formula}s سټ پوره کوي۔ نو
  يوازې دا نۀ ده چې که $\Gamma \Entails !A$ وي، تړلے
  !!{tableau} او $\Gamma \Proves !A$ هم وي؛ که د تړلې
  !!{tableau} لټون په سمه منظمه طريقه وکړو او «بشپړې»
  !!{tableau} ته
  ورسېږو، نو $\Gamma \Entails/ !A$ به هم وپېژنو او ضدمدل
  به په رښتيا جوړولے شو۔
\end{explain}

\iftag{prvBox}{
\begin{ex}
  پوهېږو چې $\Proves/ \Box(p \lor q) \lif (\Box p \lor \Box
  q)$۔ د تابلو جوړول داسې پيلېږي:
  \begin{oltableau}
    [\pFmla{\False}{\Box(p \lor q) \lif (\Box p \lor \Box q)}{1},
      just = \TAss, checked
      [\pFmla{\True}{\Box(p \lor q)}{1},
        just = {\TRule{\False}{\lif}[1]}, 
        [\pFmla{\False}{\Box p \lor \Box q}{1},
          just = {\TRule{\False}{\lif}[1]}, checked
          [\pFmla{\False}{\Box p}{1},
            just = {\TRule{\False}{\lor}[3]}, checked
            [\pFmla{\False}{\Box q}{1},
              just = {\TRule{\False}{\lor}[3]}, checked
              [\pFmla{\False}{p}{1.1}, 
                just = {\TRule{\False}{\Box}[4]}, checked
                [\pFmla{\False}{q}{1.2}, 
                  just = {\TRule{\False}{\Box}[5]}, checked
                ]
              ]
            ]
          ]
        ]
      ]
    ]
  \end{oltableau}
  دا !!{tableau} لا نۀ ده بشپړه۔ په ورپسې پړاو کښې يوازې
  هغه کرښه ګورو چې د کتنې نښه نۀ لري: په~$2$مه کرښه کښې
  مخکښ‌لرونکے !!{formula}
  $\sFmla{\True}{\Box(p \lor q)}[1]$۔ د بشپړې تابلو د
  جوړولو له مخې د $\TRule{\True}{\Box}$ قاعده بايد د څانګې
  د هر کارېدلي مخکښ لپاره، يعنې د $1.1$ او $1.2$ دواړو لپاره
  پلې شي۔ منجمده سرچينه دا لټون «تړلے» بولي، خو دلته د
  ضدمدل لپاره «بشپړه» څانګه جوړېږي:
  \begin{oltableau}
    [\pFmla{\False}{\Box(p \lor q) \lif (\Box p \lor \Box q)}{1},
      just = \TAss, checked
      [\pFmla{\True}{\Box(p \lor q)}{1},
        just = {\TRule{\False}{\lif}[1]}, 
        [\pFmla{\False}{\Box p \lor \Box q}{1},
          just = {\TRule{\False}{\lif}[1]}, checked
          [\pFmla{\False}{\Box p}{1},
            just = {\TRule{\False}{\lor}[3]}, checked
            [\pFmla{\False}{\Box q}{1},
              just = {\TRule{\False}{\lor}[3]}, checked
              [\pFmla{\False}{p}{1.1}, 
                just = {\TRule{\False}{\Box}[4]}, checked
                [\pFmla{\False}{q}{1.2}, 
                  just = {\TRule{\False}{\Box}[5]}, checked
                  [\pFmla{\True}{p \lor q}{1.1}, 
                    just = {\TRule{\True}{\Box}[2]}
                    [\pFmla{\True}{p \lor q}{1.2}, 
                      just = {\TRule{\True}{\Box}[2]}
                    ]
                  ]
                ]
              ]
            ]
          ]
        ]
      ]
    ]
  \end{oltableau}
  اوس ۲مه، ۸مه او ۹مه کرښې د کتنې نښه نۀ لري۔ خو نوے
  مخکښ نۀ دے ورزيات شوے؛ نو $\TRule{\True}{\lor}$
  په ۸مه او ۹مه کرښو او ټولو راپيدا شويو څانګو پلې کوو،
  تر هغو چې څانګه تړلې نۀ شي:
  \begin{oltableau}
    [\pFmla{\False}{\Box(p \lor q) \lif (\Box p \lor \Box q)}{1},
      just = \TAss, checked
      [\pFmla{\True}{\Box(p \lor q)}{1},
        just = {\TRule{\False}{\lif}[1]}, checked
        [\pFmla{\False}{\Box p \lor \Box q}{1},
          just = {\TRule{\False}{\lif}[1]}, checked
          [\pFmla{\False}{\Box p}{1},
            just = {\TRule{\False}{\lor}[3]}, checked
            [\pFmla{\False}{\Box q}{1},
              just = {\TRule{\False}{\lor}[3]}, checked
              [\pFmla{\False}{p}{1.1}, 
                just = {\TRule{\False}{\Box}[4]}, checked
                [\pFmla{\False}{q}{1.2}, 
                  just = {\TRule{\False}{\Box}[5]}, checked
                  [\pFmla{\True}{p \lor q}{1.1}, 
                    just = {\TRule{\True}{\Box}[2]}, checked
                    [\pFmla{\True}{p \lor q}{1.2}, 
                      just = {\TRule{\True}{\Box}[2]}, checked
                      [\pFmla{\True}{p}{1.1},
                        just = {\TRule{\True}{\lor}[8]}, checked, close
                      ]
                      [\pFmla{\True}{q}{1.1},
                        just = {\TRule{\True}{\lor}[8]}, checked
                        [\pFmla{\True}{p}{1.2},
                          just = {\TRule{\True}{\lor}[9]}, checked]
                        [\pFmla{\True}{q}{1.2},
                          just = {\TRule{\True}{\lor}[9]}, checked, close]
                      ]
                    ]
                  ]
                ]
              ]
            ]
          ]
        ]
      ]
    ]
  \end{oltableau}
  يوه پرانيستې څانګه پاتې ده او بشپړه ده۔ له هغې داسې مدل
  تعريفوو چې نړۍې يې $W = \{1, 1.1, 1.2\}$ دي، يعنې
  په پرانيستې څانګه کښې راغلي يوازيني مخکښونه؛ د لاسرسي
  اړيکه يې $R = \{\tuple{1, 1.1}, \tuple{1, 1.2}\}$ ده؛
  او ارزښت ټاکنه يې $V(p) = \{1.2\}$ ده، ځکه په ۱۱مه کرښه
  کښې $\sFmla{\True}{p}[1.2]$ شته، او $V(q) = \{1.1\}$
  ده، ځکه په ۱۰مه کرښه کښې
  $\sFmla{\True}{q}[1.1]$ شته۔ مدل په
  \olref{fig:counter-Box} کښې انځور شوے؛ کتلے شئ چې دا د
  $\Box(p \lor q) \lif (\Box p \lor \Box q)$ ضدمدل دے۔
  \begin{figure}
  \begin{center}
    \begin{tikzpicture}[modal]
      \node[world] (w1) [label={[align=right]right:\mFalse{p}\\ \mFalse{q}}]{$1$}; 
      \node[world] (w2) [label={[align=right]right:\mFalse{p}\\ \mTrue{q}},
        above left=of w1]{$1.1$}; 
      \node[world] (w3) [label={[align=right]right:\mTrue{p}\\ \mFalse{q}},
        above right=of w1] {$1.2$};
      \draw[->] (w1) to (w2);
      \draw[->] (w1) to (w3);
    \end{tikzpicture}
  \end{center}
  \caption{د $\Box(p \lor q) \lif (\Box p \lor \Box
    q)$ ضدمدل۔}
  \ollabel{fig:counter-Box}
\end{figure}
\end{ex}          
}
{
\begin{ex}
  پوهېږو چې $\Proves/ (\Diamond p \land \Diamond q) \lif \Diamond(p
  \land q)$۔ د تابلو جوړول داسې پيلېږي:
  \begin{oltableau}
    [\pFmla{\False}{(\Diamond p \land \Diamond q) \lif \Diamond(p \land q)}{1},
      just = \TAss, checked
      [\pFmla{\True}{\Diamond p \land \Diamond q}{1},
        just = {\TRule{\False}{\lif}[1]}, checked
        [\pFmla{\False}{\Diamond(p \land q)}{1},
          just = {\TRule{\False}{\lif}[1]}
          [\pFmla{\True}{\Diamond p}{1},
            just = {\TRule{\True}{\land}[2]}, checked
            [\pFmla{\True}{\Diamond q}{1},
              just = {\TRule{\True}{\land}[2]}, checked
              [\pFmla{\True}{p}{1.1}, 
                just = {\TRule{\True}{\Diamond}[4]}, checked
                [\pFmla{\True}{q}{1.2}, 
                  just = {\TRule{\True}{\Diamond}[5]}, checked
                ]
              ]
            ]
          ]
        ]
      ]
    ]
  \end{oltableau}
  دا !!{tableau} لا نۀ ده بشپړه۔ په ورپسې پړاو کښې يوازې
  هغه کرښه ګورو چې د کتنې نښه نۀ لري: په~$3$مه کرښه کښې
  مخکښ‌لرونکے !!{formula}
  $\sFmla{\False}{\Diamond(p \land q)}[1]$۔ د بشپړې تابلو
  د جوړولو له مخې د $\TRule{\False}{\Diamond}$ قاعده بايد
  د څانګې د هر کارېدلي مخکښ لپاره، يعنې د $1.1$ او $1.2$
  دواړو لپاره پلې شي۔ په منجمده سرچينه کښې دلته د ناسمې
  امکاني نښې پر ځاے رښتينې امکاني نښه او د هغې قاعده ليکل شوې؛
  د تابلو د درېيمې کرښې نښه او قاعده ناسم لوری غواړي۔
  همدارنګه د ضدمدل لپاره بشپړه، نۀ تړلې، څانګه جوړوو:
  د منجمدې سرچينې د منځنۍ ونې په لومړۍ کرښه کښې د
  استلزام دواړه خواوې سرچپه دي؛ لاندې يې د لومړۍ او
  وروستۍ ونې له اصلي استلزام سره يو شان راوړو۔
  \begin{oltableau}
    [\pFmla{\False}{(\Diamond p \land \Diamond q) \lif \Diamond(p \land q)}{1},
      just = \TAss, checked
      [\pFmla{\True}{\Diamond p \land \Diamond q}{1},
        just = {\TRule{\False}{\lif}[1]}, checked
        [\pFmla{\False}{\Diamond (p \land q)}{1},
          just = {\TRule{\False}{\lif}[1]}
          [\pFmla{\True}{\Diamond p}{1},
            just = {\TRule{\True}{\land}[2]}, checked
            [\pFmla{\True}{\Diamond q}{1},
              just = {\TRule{\True}{\land}[2]}, checked
              [\pFmla{\True}{p}{1.1}, 
                just = {\TRule{\True}{\Diamond}[4]}, checked
                [\pFmla{\True}{q}{1.2}, 
                  just = {\TRule{\True}{\Diamond}[5]}, checked
                  [\pFmla{\False}{p \land q}{1.1}, 
                    just = {\TRule{\False}{\Diamond}[3]}
                    [\pFmla{\False}{p \land q}{1.2}, 
                      just = {\TRule{\False}{\Diamond}[3]}
                    ]
                  ]
                ]
              ]
            ]
          ]
        ]
      ]
    ]
  \end{oltableau}
  اوس ۳مه، ۸مه او ۹مه کرښې د کتنې نښه نۀ لري۔ خو نوے
  مخکښ نۀ دے ورزيات شوے؛ نو $\TRule{\False}{\land}$
  په ۸مه او ۹مه کرښو او ټولو راپيدا شويو څانګو پلې کوو،
  تر هغو چې څانګه تړلې نۀ شي:
  \begin{oltableau}
    [\pFmla{\False}{(\Diamond p \land \Diamond q) \lif \Diamond(p \land q)}{1},
      just = \TAss, checked
      [\pFmla{\True}{\Diamond p \land \Diamond q}{1},
        just = {\TRule{\False}{\lif}[1]}, checked
        [\pFmla{\False}{\Diamond(p \land q)}{1},
          just = {\TRule{\False}{\lif}[1]}
          [\pFmla{\True}{\Diamond p}{1},
            just = {\TRule{\True}{\land}[2]}, checked
            [\pFmla{\True}{\Diamond q}{1},
              just = {\TRule{\True}{\land}[2]}, checked
              [\pFmla{\True}{p}{1.1}, 
                just = {\TRule{\True}{\Diamond}[4]}, checked
                [\pFmla{\True}{q}{1.2}, 
                  just = {\TRule{\True}{\Diamond}[5]}, checked
                  [\pFmla{\False}{p \land q}{1.1}, 
                    just = {\TRule{\False}{\Diamond}[3]}, checked
                    [\pFmla{\False}{p \land q}{1.2}, 
                      just = {\TRule{\False}{\Diamond}[3]}, checked
                      [\pFmla{\False}{p}{1.1},
                        just = {\TRule{\False}{\land}[8]}, checked, close
                      ]
                      [\pFmla{\False}{q}{1.1},
                        just = {\TRule{\False}{\land}[8]}, checked
                        [\pFmla{\False}{p}{1.2},
                          just = {\TRule{\False}{\land}[9]}, checked]
                        [\pFmla{\False}{q}{1.2},
                          just = {\TRule{\False}{\land}[9]}, checked, close]
                      ]
                    ]
                  ]
                ]
              ]
            ]
          ]
        ]
      ]
    ]
  \end{oltableau}
  يوه پرانيستې څانګه پاتې ده او بشپړه ده۔ له هغې داسې مدل
  تعريفوو چې نړۍې يې $W = \{1, 1.1, 1.2\}$ دي، يعنې
  په پرانيستې څانګه کښې راغلي يوازيني مخکښونه؛ د لاسرسي
  اړيکه يې $R = \{\tuple{1, 1.1}, \tuple{1, 1.2}\}$ ده؛
  او ارزښت ټاکنه يې $V(p) = \{1.1\}$ ده، ځکه په ۶مه کرښه
  کښې $\sFmla{\True}{p}[1.1]$ شته، او $V(q) = \{1.2\}$
  ده، ځکه په ۷مه کرښه کښې
  $\sFmla{\True}{q}[1.2]$ شته۔ په منجمده سرچينه کښې د
  دې وروستي شاهد مخکښ ۱٫۱ ليکل شوے، خو ونه او انځور دواړه
  ۱٫۲ ښيي۔ مدل په \olref{fig:counter-Diamond} کښې انځور
  شوے؛ کتلے شئ چې دا د
  $(\Diamond p \land \Diamond q) \lif \Diamond (p \land q)$
  ضدمدل دے۔
  \begin{figure}
  \begin{center}
    \begin{tikzpicture}[modal]
      \node[world] (w1) [label={[align=right]right:\mFalse{p}\\ \mFalse{q}}]{$1$}; 
      \node[world] (w2) [label={[align=right]right:\mTrue{p}\\ \mFalse{q}},
        above left=of w1]{$1.1$}; 
      \node[world] (w3) [label={[align=right]right:\mFalse{p}\\ \mTrue{q}},
        above right=of w1] {$1.2$};
      \draw[->] (w1) to (w2);
      \draw[->] (w1) to (w3);
    \end{tikzpicture}
  \end{center}
  \caption{د $(\Diamond p \land \Diamond q) \lif
    \Diamond (p \land q)$ ضدمدل۔}
  \ollabel{fig:counter-Diamond}
\end{figure}
\end{ex}          
}

\end{document}
